\documentclass[10pt]{beamer}
\usetheme{metropolis}

\usepackage{amsmath,amssymb}
\usepackage{mathtools}
\usepackage{xeCJK}
\setCJKmainfont{Noto Serif SC}
\usepackage{tikz-cd}

% 1. 开启标题背景色填充
\metroset{block=fill}

% theorem environments
\newtheorem{proposition}{Proposition}
\newtheorem{exercise}{Exercise}

% notation
\newcommand{\R}{\mathbb{R}}
\newcommand{\N}{\mathbb{N}}
\newcommand{\F}{\mathcal{F}}
\newcommand{\supp}{\operatorname{supp}}

\title{Submanifolds}
\subtitle{Normal Forms, Immersions, Submersions \& Regular Submanifolds}
\date{October 18, 2026}
\author{LXJ}

\begin{document}

\maketitle

%====================================================================
\section{Normal Form of an Immersion}

\begin{frame}{Immersions have a local standard form}
	\begin{theorem}[Normal form of an immersion]
		Let $f : M^{m} \longrightarrow N^{n}$ be an immersion. For every
		$p \in M$, there exists a local coordinate system $(U, \varphi)$
		about $p$ and a local coordinate system $(V, \psi)$ about
		$q = f(p)$, s.t.\ the local form of $f$
		\[
			\psi \circ f \circ \varphi^{-1} : \varphi(U) \longrightarrow \psi(V)
		\]
		is a \emph{standard embedding}, i.e.
		\[
			\psi \circ f \circ \varphi^{-1}(x^{1}, x^{2}, \dots, x^{m})
			= (x^{1}, x^{2}, \dots, x^{m}, 0, \dots, 0).
		\]
	\end{theorem}
	\begin{block}{Idea}
		An immersion looks, near each point, like the inclusion of a
		coordinate subspace $\R^{m} \hookrightarrow \R^{n}$.
	\end{block}
\end{frame}

\begin{frame}{Immersions are locally embeddings}
	\begin{corollary}
		An immersion map $f : M \longrightarrow N$ is an embedding
		\emph{locally}: around every point $p \in M$, there is a
		neighborhood $U \subset M$ such that $f\big|_{U} : U \to f(U)$
		is a homeomorphism onto its image.
	\end{corollary}
	\medskip
	\begin{block}{Remark}
		Globally an immersion need \emph{not} be an embedding --- the
		image may cross itself (e.g.\ the figure-eight curve). The local
		normal form only controls a small neighborhood of each point.
	\end{block}
\end{frame}

%====================================================================
\section{Normal Form of a Submersion}

\begin{frame}{Submersions have a local standard form}
	\begin{theorem}[Normal form of a submersion]
		Let $f : M^{m} \longrightarrow N^{n}$ be a submersion ($m \geq n$).
		Then $\forall p \in M$, there is a local coordinate system
		$(U, \varphi)$ of $p$ and $(V, \psi)$ of $q = f(p)$, s.t.\ the
		local form
		\[
			\psi \circ f \circ \varphi^{-1} : \varphi(U) \longrightarrow \psi(V)
		\]
		is a \emph{standard submersion}, i.e.
		\[
			\psi \circ f \circ \varphi^{-1}(x_{1}, \dots, x_{m})
			= (x_{1}, \dots, x_{n}).
		\]
	\end{theorem}
	\begin{block}{Idea}
		A submersion looks, near each point, like the projection
		$\R^{m} \twoheadrightarrow \R^{n}$ onto the first $n$ coordinates.
	\end{block}
\end{frame}

\begin{frame}{Submersions are open maps}
	\begin{corollary}
		A submersion $f : M \longrightarrow N$ is an \emph{open map},
		i.e.\ it maps open sets to open sets.
	\end{corollary}
	\begin{proof}
		Let $W \subset M$ be open and $p \in W$. In the normal-form
		coordinates above, $f$ is a coordinate projection, which is open
		on $\varphi(U \cap W)$; hence $f(U \cap W)$ is a neighborhood of
		$f(p)$ contained in $f(W)$. Since $p \in W$ is arbitrary, $f(W)$
		is open.
	\end{proof}
	\medskip
	\begin{block}{Remark}
		In particular, a \emph{surjective} submersion $f : M \to N$ is a
		quotient map: $V \subset N$ is open iff $f^{-1}(V) \subset M$ is
		open.
	\end{block}
\end{frame}

%====================================================================
\section{Submanifolds}

\begin{frame}{Two notions of submanifolds}
	\begin{definition}[Submanifolds]
		Let $M, N$ be $C^{\infty}$-manifolds with $M \subset N$ as a set.
		If the inclusion map $i : M \hookrightarrow N$ is an immersion,
		then $M$ is an \emph{immersion submanifold} (or submfd). If
		$i : M \hookrightarrow N$ is an embedding (i.e.\ $M \subset N$
		as a topological subspace), then $M$ is called a
		\emph{regular submanifold}.
	\end{definition}
	\medskip
	\begin{lemma}
		Let $M$ and $N$ be $C^{\infty}$-manifolds. If $f : M
			\longrightarrow N$ is a one-to-one immersion, then $f(M) \subset N$
		is an immersion submanifold (with the topological structure
		inherited from $M$). If $f : M \longrightarrow N$ is an embedding,
		then $f(M) \subset N$ is a regular submanifold.
	\end{lemma}
\end{frame}

\begin{frame}{Immersion submanifold vs.\ regular submanifold}
	\begin{block}{Terminology (following the lemma)}
		When $f : M \longrightarrow N$ is a one-to-one immersion, we call
		$M$ (or its image) an \emph{immersion submanifold}. If $f$ is an
		embedding, we call $M$ a \emph{regular submanifold}.
	\end{block}
	\medskip
	\begin{alertblock}{Key distinction}
		\begin{itemize}
			\item \textbf{Immersion submanifold:} the topology on $M$ may
			      be \emph{finer} than the subspace topology inherited from
			      $N$; ``short'' pieces of $M$ may be far apart in $N$.
			\item \textbf{Regular submanifold:} $M$ carries exactly the
			      subspace topology; locally $M$ looks like a
			      \emph{slice} of $N$ --- see the structure theorem below.
		\end{itemize}
	\end{alertblock}
\end{frame}

%====================================================================
\section{The Structure of Regular Submanifolds}

\begin{frame}{Regular submanifolds are locally slices}
	\begin{theorem}[The structure of regular submanifolds]
		$M$ is a regular submanifold of $N$ $\Longleftrightarrow$ $M$ is
		a sub-topological space of $N$, and $\forall p \in M$, there is a
		local coordinate system $(U, x_{1}, \dots, x_{n})$ of $p$ in $N$,
		s.t.
		\[
			M \cap U
			= \bigl\{ q \in U \bigm|
			x_{i}(q) = 0,\ m+1 \leq i \leq n \bigr\}.
		\]
	\end{theorem}
	\begin{block}{Remark}
		$\{ x_{i} \}_{m+1 \leq i \leq n}$ are called \emph{slice
			coordinates}: locally $M$ is cut out by setting the last
		$n - m$ coordinates to zero, and $(x_{1}, \dots, x_{m})\big|_{M}$
		give a chart of $M$.
	\end{block}
\end{frame}

\begin{frame}{The constant rank theorem}
	\begin{theorem}[Constant rank map]
		Let $f : M^{m} \longrightarrow N^{n}$ be a $C^{\infty}$ map with
		\[
			\operatorname{rank}_{p} f \equiv \mathrm{const} = l
			\qquad (\forall p \in M).
		\]
		Then $\forall q \in N$, the level set
		\[
			f^{-1}(q) = \{ p \in M \mid f(p) = q \}
		\]
		is either $\varnothing$, or a regular submanifold of $M$ of
		dimension $(m - l)$.
	\end{theorem}
	\medskip
	\begin{block}{Special cases}
		\begin{itemize}
			\item $l = n$ (submersion): $f^{-1}(q)$ is a regular
			      submanifold of dimension $m - n$ --- the preimage
			      theorem.
			\item $l = m$ (immersion, $m \leq n$): $f^{-1}(q)$ is a
			      $0$-dimensional submanifold (discrete points).
		\end{itemize}
	\end{block}
\end{frame}

%====================================================================
\section{Proof of the Structure Theorem}

\begin{frame}{Setup of the proof ($\Rightarrow$)}
	Let $M$ be a regular submanifold of $N$. The inclusion map
	\[
		i : M \longrightarrow N
	\]
	is an embedding, where $M$ is endowed with the subspace topology.
	The differential
	\[
		i_{*, p} : T_{p} M \longrightarrow T_{p} N
	\]
	has $\operatorname{rank} i_{*, p} = m$ ($m = \dim M$). Let $p \in M$;
	take two local coordinate systems about $p$:
	\[
		(U, \varphi = (x_{1}, \dots, x_{m})) \text{ on } M,
		\qquad
		(V, \psi = (y_{1}, \dots, y_{n})) \text{ on } N.
	\]
\end{frame}

\begin{frame}{Full rank of the Jacobian}
	Then the Jacobian
	\[
		\left.
		\frac{\partial \bigl( y_{i} \circ \varphi^{-1} \bigr)}
		{\partial x_{j}}
		\right|_{\substack{1 \leq i \leq n \\ 1 \leq j \leq m}}
		\bigl( \varphi(p) \bigr)
		\quad \text{has full rank } m.
	\]
	WLOG (relabeling the $y_{i}$'s), we assume the $m \times m$
	submatrix
	\[
		\left.
		\frac{\partial \bigl( y_{i} \circ \varphi^{-1} \bigr)}
		{\partial x_{j}}
		\right|_{\substack{1 \leq i \leq m \\ 1 \leq j \leq m}}
		\bigl( \varphi(p) \bigr)
		\quad \text{is non-singular.}
	\]
\end{frame}

\begin{frame}{Extending to a local diffeomorphism}
	Consider the new map $\Phi : \varphi(U) \times \R^{n-m}
		\longrightarrow \psi(V)$,
	\begin{align*}
		 & (x_{1}, \dots, x_{m}; x_{m+1}, \dots, x_{n}) \\
		 & \qquad \longmapsto
		\bigl(
		y_{1} \circ \varphi^{-1},\, \dots,\, y_{m} \circ \varphi^{-1},\
		y_{m+1} \circ \varphi^{-1} + x_{m+1},\, \dots,\,
		y_{n} \circ \varphi^{-1} + x_{n}
		\bigr).
	\end{align*}
	Then
	\[
		D \Phi \Big|_{(\varphi(p),\, 0)}
		=
		\begin{pmatrix}
			\left( \dfrac{\partial \, y_{i} \circ \varphi^{-1}}
			{\partial x_{j}} \right)_{\! m \times m} & 0               \\[8pt]
			*                                                   & I_{\, n-m}
		\end{pmatrix}
	\]
	is non-singular, so $\Phi$ is a \emph{local diffeomorphism} about
	$(\varphi(p), 0)$.
\end{frame}

\begin{frame}{Conclusion of the proof}
	Consider
	\[
		\Phi^{-1} \circ \psi \circ i \circ \varphi^{-1}
		(x_{1}, \dots, x_{m})
		= (x_{1}, \dots, x_{m},\, 0, \dots, 0).
	\]
	Indeed, by the definition of $\Phi$, on the slice
	$(x_{m+1}, \dots, x_{n}) = (0, \dots, 0)$ we have
	$\Phi \circ (\psi \circ i \circ \varphi^{-1}) (x_{1}, \dots, x_{m})
		= (x_{1}, \dots, x_{m}, 0, \dots, 0)$.
	Hence the new coordinates on $N$
	\[
		\widetilde{x}_{i} \triangleq (\Phi^{-1} \circ \psi)_{i},
		\qquad i = 1, \dots, n,
	\]
	satisfy
	\[
		M \cap U
		= \bigl\{ q \in U \bigm|
		\widetilde{x}_{i}(q) = 0,\ m+1 \leq i \leq n \bigr\}.
		\qed
	\]
\end{frame}

\end{document}
